\chapter{萬章章句上}
\kaishi*

\jiazhu[1]{萬章者，萬姓章名，孟子弟子也。萬章問舜孝，猶《論語》顔淵問仁，因以題篇。}

\section{萬章問舜孝}

萬章問曰：「舜往于田，號泣于旻天，何爲其號泣也？」\jiazhu[1]{問舜往至于田，何爲號泣也？謂耕於歷山之時。}

孟子曰：「怨慕也。」\jiazhu[1]{言舜自怨遭父母見惡之厄，而思慕也。}

萬章曰：「父母愛之，喜而不忘；父母惡之，勞而不怨。然則舜怨乎？」\jiazhu[1]{言孝法當不怨，如是舜何故怨？}

曰：「長息問於公明高曰：『舜往于田，則吾既得聞命矣；號泣于旻天、于父母，則吾不知也。』公明高曰：『是非爾所知也。』」\jiazhu[1]{長息，公明高弟子；公明高，曾子弟子。旻天，秋也，憂陰氣也，故訴于旻天。高非息之問不得其義，故曰非爾所知。}

「夫公明高以孝子之心爲不若是恝。」\jiazhu[1]{恝，無愁之貌。孟子以萬章之問難自距之，故爲言高息之相對如此。夫公明高以爲孝子不得意於父母，自當怨悲，豈可恝恝然無憂哉？因爲萬章具陳其意。}

「我竭力耕田，共爲子職而已矣。父母之不我愛，於我何哉？」\jiazhu[1]{我共人子之事，而父母不我愛，於我之身獨有何罪哉？自求責於己而悲感焉。}

「帝使其子九男二女，百官、牛羊、倉廩備，以事舜於畎畝之中。」\jiazhu[1]{帝，堯也。堯使九子事舜以爲師，以二女妻舜，百官致牛羊，倉廩致粟米之餼，備具饋禮以奉事舜於畎畝之中。由是遂賜舜以倉廩牛羊，使得自有之。《堯典》曰「釐降二女」，不見九男。孟子時，尚書凡百二十篇，逸書有《舜典》之敘，亡失其文。孟子諸所言舜事，皆《堯典》及逸書所載。獨丹朱以胤嗣之子臣下以距堯求禪，其餘八庶無事故不見於《堯典》，猶晉獻公之子九人，五人以事見於《春秋》，其餘四子亦不復見。}

「天下之士多就之者，帝將胥天下而遷之焉。爲不順於父母，如窮人無所歸。」\jiazhu[1]{天下之善士多就舜而悅之。胥，須也。堯須天下悉治，將遷位而禪之。順，愛也。爲不愛於父母，其爲憂愁若困窮之人無所歸往也。}

「天下之士悅之，人之所欲也，而不足以解憂；好色，人之所欲，妻帝之二女，而不足以解憂；富，人之所欲，富有天下，而不足以解憂；貴，人之所欲，貴爲天子，而不足以解憂。人悅之、好色、富貴，無足以解憂者，惟順於父母，可以解憂。」\jiazhu[1]{言爲人所悅，將見禪爲天子，皆不足以解憂，獨見愛於父母爲可以解己之憂。}

「人少則慕父母；知好色則慕少艾；有妻子則慕妻子；仕則慕君，不得於君則熱中。」\jiazhu[1]{慕，思慕也。人少，年少也。艾，美好也。不得於君，失意於君也。熱中，心熱恐懼也。是乃人之情。}

「大孝終身慕父母。五十而慕者，予於大舜見之矣。」\jiazhu[1]{大孝之人，終身慕父母。若老萊子七十而慕，衣五彩之衣，爲嬰兒匍匐於父母前也。我於大舜見五十而尚慕父母。《書》曰：「舜生三十徵庸，三十在位。」在位時尚慕，故言五十也。}\par \jiazhu[1]{章指言：夫孝者，百行之本，無物以先之。雖富有天下，而不能取悅於其父母，莫有可也。孝道明著，則六合歸仁矣。}

\section{萬章問舜不告而娶}

萬章問曰：「《詩》云：『娶妻如之何？必告父母。』信斯言也，宜莫如舜。舜之不告而娶，何也？」\jiazhu[1]{《詩·齊風·南山》之篇。言娶妻之禮，必告父母。舜合信此詩之言，何爲違禮不告而娶也？}

孟子曰：「告則不得娶。男女居室，人之大倫也。如告，則廢人之大倫，以懟父母，是以不告也。」\jiazhu[1]{舜父頑母嚚，常欲害舜。告則不聽其娶，是廢人之大倫，以怨懟於父母也。}

萬章曰：「舜之不告而娶，則吾既得聞命矣。帝之妻舜而不告，何也？」\jiazhu[1]{禮，娶須五禮，父母亢答以辭，是相告也。帝謂堯也，何不告舜父母也？}

曰：「帝亦知告焉則不得妻也。」\jiazhu[1]{帝堯知舜大孝，父母止之，舜不敢違，則不得妻之，故亦不告。}

萬章曰：「父母使舜完廩，捐階，瞽瞍焚廩。使浚井，出，從而揜之。」\jiazhu[1]{完，治廩倉。階，梯也。使舜登廩屋而捐去其階，焚燒其廩也。一說：旋階，舜即旋從階下，瞽瞍不知其已下，故焚廩也。使舜浚井，舜入而即出，瞽瞍不知其已出，從而蓋其井，以爲死矣。}

「象曰：『謨蓋都君，咸我績。』」\jiazhu[1]{象，舜異母弟。謨，謀；蓋，覆；都，於也。君，舜也。舜有牛羊倉廩之奉，故謂之君。咸，皆；績，功也。象言謀覆於君而殺之者，皆我之功，欲與父母分舜之有，取其善者，故引其功也。}

「牛羊父母，倉廩父母。」\jiazhu[1]{欲以牛羊倉廩與其父母。}

「干戈朕，琴朕，弤朕，二嫂使治朕棲。」\jiazhu[1]{干，楯；戈，戟也。琴，舜所彈五弦琴也。弤，雕弓也。天子曰雕弓，堯禪舜天下，故賜之雕弓也。棲，牀也。二嫂，娥皇、女英，使治牀，欲以爲妻也。}

「象往入舜宮，舜在牀琴。象曰：『鬱陶思君爾。』忸怩。」\jiazhu[1]{象見舜生在牀鼓琴，愕然，反辭曰：「我鬱陶思君，故來爾。」辭也。忸怩而慙，是其情也。}

「舜曰：『惟茲臣庶，汝其于予治。』」\jiazhu[1]{茲，此也。象素憎舜，不至其宮也，故舜見來而喜曰：「惟念此臣衆，汝故助我治事。」}

「不識舜不知象之將殺已與？」\jiazhu[1]{萬章言我不知舜不知象之將殺之與？何爲好言順辭以答象也。}

曰：「奚而不知也？象憂亦憂，象喜亦喜。」\jiazhu[1]{奚，何也。孟子曰：舜何爲不知象惡已也？仁人愛其弟，憂喜隨之。象方言思君，故以順辭答之。}

曰：「然則舜僞喜者與？」\jiazhu[1]{僞，詐也。萬章言如是，則爲舜行至誠而詐喜以悅人矣。}

曰：「否。昔者有饋生魚於鄭子產，子產使校人畜之池。校人烹之，反命曰：『始舍之，圉圉焉；少則洋洋焉；攸然而逝。』子產曰：『得其所哉！得其所哉！』」\jiazhu[1]{孟子言否，云舜不詐喜也。因爲說子產以喻之。子產，鄭子國之子公孫僑，大賢人也。校人，主池沼小吏也。圉圉，魚在水羸劣之貌。洋洋，舒緩搖尾之貌。攸然，迅走水趣深處也。故曰得其所哉，重言之，嘉得魚之志也。}

「校人出，曰：『孰謂子產智？予既烹而食之，曰：得其所哉，得其所哉。』故君子可欺以其方，難罔以非其道。彼以愛兄之道來，故誠信而喜之，奚僞焉？」\jiazhu[1]{方，類也。君子可以事類欺，故子產不知校人之食其魚。象以其愛兄之言來向舜，是亦其類也，故誠信之而喜，何爲僞喜也？}\par \jiazhu[1]{章指言：仁聖所存者大，舍小從大，達權之義也。不告而娶，守正道也。}

\section{萬章問象封有庳}

萬章問曰：「象日以殺舜爲事，立爲天子則放之，何也？」\jiazhu[1]{怪舜放之何故？}

孟子曰：「封之也。或曰放焉。」\jiazhu[1]{舜封象于有庳，或有人以爲放之。}

萬章曰：「舜流共工于幽州，放驩兜于崇山，殺三苗于三危，殛鯀于羽山，四罪而天下咸服，誅不仁也。象至不仁，封之有庳。有庳之人奚罪焉？仁人固如是乎？在他人則誅之，在弟則封之。」\jiazhu[1]{舜誅四佞以其惡也。象惡亦甚而封之，仁人用心當如是乎？罪在他人當誅之，在弟則封之。}

曰：「仁人之於弟也，不藏怒焉，不宿怨焉，親愛之而已矣。親之欲其貴也，愛之欲其富也。封之有庳，富貴之也。身爲天子，弟爲匹夫，可謂親愛之乎？」\jiazhu[1]{孟子言仁人於弟，不問善惡，親愛之而已。封者欲使富貴耳。身爲天子，弟雖不仁，豈可使爲匹夫也？}

「敢問或曰放者，何謂也？」\jiazhu[1]{萬章問放之意。}

曰：「象不得有爲於其國，天子使吏治其國而納其貢稅焉，故謂之放。豈得暴彼民哉？」\jiazhu[1]{象不得施教於其國，天子使吏代其治而納貢賦與之，比諸見放也。有庳雖不得賢君，象亦不侵其民也。}

「雖然，欲常常而見之，故源源而來，不及貢，以政接于有庳。」\jiazhu[1]{雖不使象得豫政事，舜以兄弟之恩，欲常常見之無已，故源源而來，如流水之與源通。不及貢者，不待朝貢諸侯常禮乃來也。其間歲歲自至京師，謂若天子以政事接見有庳之君者，實親親之恩也。}

「此之謂也。」\jiazhu[1]{此常常已下，皆尚書逸篇之辭。孟子以告萬章，言此乃象之謂也。}\par \jiazhu[1]{章指言：懇誠于內者，則外發於事，仁人之心也。象爲無道極矣，友于之性，忘其悖逆，況其仁賢乎？}

\section{咸丘蒙問舜臣父}

咸丘蒙問曰：「語云：『盛德之士，君不得而臣，父不得而子。』舜南面而立，堯帥諸侯北面而朝之，瞽瞍亦北面而朝之。舜見瞽瞍，其容有蹙。孔子曰：『於斯時也，天下殆哉，岌岌乎！』不識此語誠然乎哉？」\jiazhu[1]{咸丘蒙，孟子弟子。語者，諺語也。言盛德之士，君不敢臣，父不敢子。堯與瞽瞍皆臣事舜，其容有蹙踖不自安也。孔子以爲君父爲臣，岌岌乎不安貌也，故曰殆哉。不知此語實然乎？}

孟子曰：「否。」\jiazhu[1]{言不然也。}

「此非君子之言，齊東野人之語也。」\jiazhu[1]{東野，東作田野之人所言耳。咸丘蒙，齊人也，故問齊野人之言。《書》曰「平秩東作」，謂治農事也。}

「堯老而舜攝也。《堯典》曰：『二十有八載，放勛乃徂落，百姓如喪考妣，三年，四海遏密八音。』」\jiazhu[1]{孟子言舜攝行事耳，未爲天子也。放勛，堯名。徂落，死也。如喪考妣，思之如父母也。遏，止也；密，無聲也。八音不作，哀思甚也。}

「孔子曰：『天無二日，民無二王。』舜既爲天子矣，又帥天下諸侯以爲堯三年喪，是二天子矣。」\jiazhu[1]{日一，王一，言不得竝也。}

咸丘蒙曰：「舜之不臣堯，則吾既得聞命矣。」\jiazhu[1]{不以堯爲臣也。}

「《詩》云：『普天之下，莫非王土；率土之濱，莫非王臣。』而舜既爲天子矣，敢問瞽瞍之非臣，如何？」\jiazhu[1]{《詩·小雅·北山》之篇。普，遍；率，循也。遍天下，循土之濱，無有非王者之臣，而曰瞽瞍非臣，如何也？}

曰：「是詩也，非是之謂也。勞於王事而不得養父母也。曰：『此莫非王事，我獨賢勞也。』」\jiazhu[1]{孟子言此詩非舜臣父之謂也。詩言皆王臣也，何爲獨使我以賢才而勞苦，不得養父母乎？是以怨也。}

「故說詩者，不以文害辭，不以辭害志。以意逆志，是爲得之。如以辭而已矣，《雲漢》之詩曰：『周餘黎民，靡有孑遺。』信斯言也，是周無遺民也。」\jiazhu[1]{文，詩之文章所引以興事也。辭，詩人所歌詠之辭。志，詩人志所欲之事。意，學者之心意也。孟子言說詩者當本之，不可以文害其辭，文不顯乃反顯也；不可以辭害其志，辭曰「周餘黎民，靡有孑遺」，志在憂旱災民無孑然遺脫不遭旱災者，非無民也。人情不遠，以己之意逆詩人之志，是爲得其實矣。王者有所不臣，不可謂皆爲王臣，謂舜臣父也。}

「孝子之至，莫大乎尊親；尊親之至，莫大乎以天下養。爲天子父，尊之至也；以天下養，養之至也。」\jiazhu[1]{尊之至，瞽瞍爲天子父；養之至，舜以天下之富奉養其親，至極也。}

「《詩》曰：『永言孝思，孝思惟則。』此之謂也。」\jiazhu[1]{《詩·大雅·下武》之篇，周武王所以長言孝道，欲以爲天下法則。此舜之謂也。}

「《書》曰：『祗載見瞽瞍，夔夔齋慄，瞽瞍亦允若。』是爲父不得而子也。」\jiazhu[1]{《書》，尚書逸篇。祗，敬；載，事也。夔夔齋慄，敬慎戰懼貌。舜既爲天子，敬事嚴父，戰慄以見瞽瞍，瞍亦信知舜之大孝。若是爲父不得而子也，以是解咸丘蒙之疑。}\par \jiazhu[1]{章指言：孝莫大於嚴父而尊之矣，行莫過於蒸蒸執子之政也。此聖人之軌道，無有加焉。}

\section{萬章問堯與舜天下}

萬章問曰：「堯以天下與舜，有諸？」\jiazhu[1]{欲知堯實以天下與舜否？}

孟子曰：「否。天子不能以天下與人。」\jiazhu[1]{堯不與之。當與天意合之，非天命者，天子不能違天命也。堯曰「咨爾舜，天之曆數在爾躬」是也。}

「然則舜有天下也，孰與之？」\jiazhu[1]{萬章言誰與之也。}

曰：「天與之。」\jiazhu[1]{孟子言天與之。}

「天與之者，諄諄然命之乎？」\jiazhu[1]{萬章言天有聲音命與之乎？}

曰：「否。天不言，以行與事示之而已矣。」\jiazhu[1]{孟子曰：天不言語，但以其人之所行善惡，又以其事從而示天下也。}

曰：「以行與事示之者，如之何？」\jiazhu[1]{萬章欲知示之之意。}

曰：「天子能薦人於天，不能使天與之天下；諸侯能薦人於天子，不能使天子與之諸侯；大夫能薦人於諸侯，不能使諸侯與之大夫。昔者堯薦舜於天而天受之，暴之於民而民受之。故曰：天不言，以行與事示之而已矣。」\jiazhu[1]{孟子言下能薦人於上，不能令上必用之。舜，天人所受，故得天下也。}

曰：「敢問薦之於天而天受之，暴之於民而民受之，如何？」\jiazhu[1]{萬章言天人受之，其事云何？}

曰：「使之主祭而百神享之，是天受之；使之主事而事治，百姓安之，是民受之也。天與之，人與之，故曰天子不能以天下與人。」\jiazhu[1]{百神享之，祭祀得福也。百姓安之，民皆謳歌其德也。}

「舜相堯二十有八載，非人之所能爲也，天也。」\jiazhu[1]{二十八年之久，非人爲也，天與之也。}

「堯崩，三年之喪畢，舜避堯之子於南河之南。天下諸侯朝覲者，不之堯之子而之舜；訟獄者，不之堯之子而之舜；謳歌者，不謳歌堯之子而謳歌舜。故曰天也。夫然後之中國，踐天子位焉。而居堯之宮，逼堯之子，是篡也，非天與也。」\jiazhu[1]{南河之南，遠地南夷也。故言然後之中國。堯子，胤子丹朱。訟獄，獄不決其罪，故訟之。謳歌，舜德也。}

「《太誓》曰：『天視自我民視，天聽自我民聽。』此之謂也。」\jiazhu[1]{《太誓》，尚書篇名。自，從也。言天之視聽從人所欲也。}\par \jiazhu[1]{章指言：德合於天，則天爵歸之；行歸於仁，則天下與之。天命不常，此之謂也。}

\section{萬章問禹德衰}

萬章問曰：「人有言：『至於禹而德衰，不傳於賢而傳於子。』有諸？」\jiazhu[1]{問禹之德衰，不傳於賢而自傳於子，有之否？}

孟子曰：「否，不然也。」\jiazhu[1]{否，不也。不如人所言。}

「天與賢，則與賢；天與子，則與子。」\jiazhu[1]{言隨天也。}

「昔者舜薦禹於天，十有七年。舜崩，三年之喪畢，禹避舜之子於陽城。天下之民從之，若堯崩之後不從堯之子而從舜也。禹薦益於天，七年。禹崩，三年之喪畢，益避禹之子於箕山之陰。朝覲訟獄者不之益而之啓，曰：『吾君之子也。』謳歌者不謳歌益而謳歌啓，曰：『吾君之子也。』」\jiazhu[1]{舜薦禹，禹薦益，同也。以啓之賢，故天下歸之。益又未久故也。陽城、箕山之陰，皆嵩山下深谷之中以藏處也。}

「丹朱之不肖，舜之子亦不肖。舜之相堯、禹之相舜也，歷年多，施澤於民久。啓賢，能敬承繼禹之道。益之相禹也，歷年少，施澤於民未久。」\jiazhu[1]{舜薦禹、禹薦益同也。以啓之賢，故天下歸之。益又未久故也。}

「舜、禹、益相去久遠，其子之賢不肖，皆天也，非人之所能爲也。莫之爲而爲者，天也；莫之致而至者，命也。」\jiazhu[1]{莫，無也。人無所欲爲而橫爲之者，天使爲也。人無欲致此事而此事自至者，是其命祿也。}

「匹夫而有天下者，德必若舜、禹，而又有天子薦之者，故仲尼不有天下。」\jiazhu[1]{仲尼無天子之薦，故不得有天下。}

「繼世以有天下。」\jiazhu[1]{繼世之君雖無仲尼之德，襲父之位，非匹夫，故得有天下也。}

「天之所廢，必若桀、紂者也，故益、伊尹、周公不有天下。」\jiazhu[1]{益值啓之賢，伊尹值太甲能改過，周公值成王有德，不遭桀、紂，故以匹夫而不有天下。}

「伊尹相湯以王於天下。湯崩，大丁未立，外丙二年，仲壬四年。太甲顛覆湯之典刑，伊尹放之於桐。三年，太甲悔過，自怨自艾，於桐處仁遷義，三年，以聽伊尹之訓己也，復歸于亳。」\jiazhu[1]{大丁，湯之大子，未立而薨。外丙立二年，仲壬立四年，皆大丁之弟也。太甲，大丁子也。伊尹以其顛覆典刑，放之於桐邑。處，居也；遷，徙也。居仁徙義，自怨其惡行。艾，治也，治而改過，以聽伊尹之教訓己，故復得歸之於亳，反天子位也。}

「周公之不有天下，猶益之於夏、伊尹之於殷也。」\jiazhu[1]{周公與益、伊尹雖有聖賢之德，不遭者時然。}

「孔子曰：『唐、虞禪，夏后、殷、周繼，其義一也。』」\jiazhu[1]{孔子言禪、繼其義一也。}\par \jiazhu[1]{章指言：施仁於民，則四海宅心；守正不足，則聖位莫繼。丹朱、商均是也。是以聖人孜孜於仁德也。}

\section{萬章問伊尹以割烹要湯}

萬章問曰：「人有言：『伊尹以割烹要湯。』有諸？」\jiazhu[1]{人言伊尹負鼎俎而干湯，有之否？}

孟子曰：「否，不然。」\jiazhu[1]{否，不是也。}

「伊尹耕於有莘之野，而樂堯舜之道焉。非其義也，非其道也，祿之以天下弗顧也；繫馬千駟，弗視也。非其義也，非其道也，一介不以與人，一介不以取諸人。」\jiazhu[1]{有莘，國名。伊尹初隱之時，耕於有莘之國，樂仁義之道。非仁義之道者，雖以天下之祿加之，不一顧而覦也。千駟，四千匹也，雖多不一眄視也。一介草不以與人，亦不以取於人也。}

「湯使人以幣聘之，囂囂然曰：『我何以湯之聘幣爲哉？我豈若處畎畝之中，由是以樂堯舜之道哉？』」\jiazhu[1]{湯聞其賢，以玄纁之幣帛往聘之。囂囂，自得之志，無欲之貌也。曰豈若居畎畝之中而無憂哉？樂我堯舜仁義之道。}

「湯三使往聘之。既而幡然改曰：『與我處畎畝之中，由是以樂堯舜之道，吾豈若使是君爲堯舜之君哉？吾豈若使是民爲堯舜之民哉？吾豈若於吾身親見之哉？』」\jiazhu[1]{幡，反也。三聘既至，而後幡然改本之計，欲就湯聘以行其道，使君爲堯舜之君，使民爲堯舜之民。}

「天之生此民也，使先知覺後知，使先覺覺後覺也。予，天民之先覺者也。予將以斯道覺斯民也。非予覺之，而誰也？」\jiazhu[1]{覺，悟也。天欲使先知之人悟後知之人。我先悟覺者也，我欲以此仁義之道覺悟此未知之民，非我悟之，將誰教乎？}

「思天下之民，匹夫匹婦有不被堯舜之澤者，若己推而內之溝中。其自任以天下之重如此，故就湯而說之以伐夏救民。」\jiazhu[1]{伊尹思念不以仁義之道化民者，如己推排內之溝壑中也。自任其重如此，故就湯說之伐夏桀，救民之厄也。}

「吾未聞枉己而正人者也，況辱己以正天下者乎？」\jiazhu[1]{枉己者尚不能以正人，況於辱己之身而有正天下者也？}

「聖人之行不同也，或遠或近，或去或不去，歸潔其身而已矣。」\jiazhu[1]{不同，謂所由不同，大要當同歸，但殊塗耳。或遠者，處身遠也；或近者，仕者近君也；或去者，不屑就也；或不去者，云焉能浼我也？歸於身潔，不污己而已。}

「吾聞其以堯舜之道要湯，未聞以割烹也。」\jiazhu[1]{我聞伊尹以仁義干湯，致湯爲王，不聞以割烹牛羊爲道。}

「《伊訓》曰：『天誅造攻自牧宮，朕載自亳。』」\jiazhu[1]{《伊訓》，尚書逸篇名。牧宮，桀宮。朕，我也，謂湯也。載，始也。亳，殷都也。言意欲誅伐桀，造作可攻討之罪者，從牧宮桀起，自取之也。湯曰我始與伊尹謀之於亳，遂順天而誅也。}\par \jiazhu[1]{章指言：賢達之理世務也，推正以濟時物，守己直行，不枉道而取容，期於益治而已矣。}

\section{萬章問孔子主癰疽}

萬章問曰：「或謂孔子於衛主癰疽，於齊主侍人瘠環，有諸乎？」\jiazhu[1]{有人以孔子爲然。癰疽，癰疽之醫也。瘠，姓；環，名，侍人也。衛君、齊君之所近狎人。}

孟子曰：「否，不然也。好事者爲之也。」\jiazhu[1]{否，不也。不如是也。好事毀人德行者爲之辭也。}

「於衛主顔讎由。彌子之妻與子路之妻，兄弟也。彌子謂子路曰：『孔子主我，衛卿可得也。』子路以告，孔子曰：『有命。』孔子進以禮，退以義，得之不得曰有命。而主癰疽與侍人瘠環，是無義無命也。」\jiazhu[1]{顔讎由，衛賢大夫，孔子以爲主。彌子，彌子瑕也，因子路欲爲孔子主。孔子知彌子幸於靈公，不以正道，故不納之而歸於命也。孔子進以禮，退應義，必曰有天命也。若主此二人，是爲無義無命也。}

「孔子不悅於魯、衛，遭宋桓司馬將要而殺之，微服而過宋。是時孔子當阨，主司城貞子，爲陳侯周臣。」\jiazhu[1]{孔子以道不合，不見悅魯、衛之君而去。適諸侯，遭宋桓魋之故，乃變更微服而過宋。司城貞子，宋卿也，雖非大賢，亦無諂惡之罪，故謚爲貞子。陳侯周，陳懷公子也，爲楚所滅，故無諡，但曰陳侯周。是時孔子遭阨難，不暇擇大賢臣而主貞子，爲陳侯周臣也。於衛、齊無阨難，何爲主癰疽瘠環也？}

「吾聞觀近臣，以其所爲主；觀遠臣，以其所主。若孔子主癰疽與侍人瘠環，何以爲孔子？」\jiazhu[1]{近臣當爲遠方來賢者爲主。遠臣自遠而至，當主於在朝之臣賢者。若孔子主於卑幸之臣，是爲凡人耳，何謂孔子得見稱爲聖人？}\par \jiazhu[1]{章指言：君子大居正，以禮進退，屈伸達節，不違貞信，故孟子辯之，正其大義也。}

\section{萬章問百里奚自鬻}

萬章問曰：「或曰：『百里奚自鬻於秦養牲者，五羊之皮，食牛以要秦繆公。』信乎？」\jiazhu[1]{人言百里奚自賣五羖羊皮，爲人養牛，以是而要繆公之相，實然不？}

孟子曰：「否，不然。好事者爲之也。」\jiazhu[1]{好事毀敗人之德行者，爲之設此言。}

「百里奚，虞人也。晉人以垂棘之璧與屈產之乘，假道於虞以伐虢。宮之奇諫。」\jiazhu[1]{垂棘，美玉所出地名。屈產，地良馬所生。乘，四馬也。皆晉國之所寶。宮之奇，虞之賢臣，諫不欲令虞公受璧馬假晉道。}

「百里奚不諫，知虞公之不可諫而去之秦，年已七十矣。曾不知以食牛干秦繆公之爲汙也，可謂智乎？不可諫而不諫，可謂不智乎？知虞公之將亡而先去之，不可謂不智也。時舉於秦，知繆公之可與有行也而相之，可謂不智乎？相秦而顯其君於天下，可傳於後世，不賢而能之乎？」\jiazhu[1]{百里奚知虞公之不可諫而去之秦，年七十而不知食牛干人君之爲汙，是爲不智也。欲言其不智，下有三智。知食牛干秦爲不然也。卒相秦，顯其君，不賢之人豈能如是？言其實賢也。}

「自鬻以成其君，鄉黨自好者不爲，而謂賢者爲之乎？」\jiazhu[1]{人自鬻於汙辱而以傅相成立其君。鄉黨邑里自喜好名者尚不肯爲也，況賢人肯辱身而爲之乎？}\par \jiazhu[1]{章指言：君子時行則行，時舍則舍，故能顯君明道，不爲苟合而違正也。}
